% Generated draft: requires mathematical and source review.
\documentclass[11pt]{article}
\usepackage[margin=25mm]{geometry}
\usepackage{fontspec}
\setmainfont{Latin Modern Roman}
\usepackage{amsmath,amssymb,mathtools,unicode-math}
\setmathfont{Latin Modern Math}
\usepackage{enumitem,needspace,xurl,hyperref}
\setlength{\parindent}{0pt}
\setlength{\parskip}{5pt}
\setlength{\emergencystretch}{3em}
\newcommand{\N}{\mathbb{N}}
\newcommand{\Z}{\mathbb{Z}}
\newcommand{\Q}{\mathbb{Q}}
\newcommand{\R}{\mathbb{R}}
\newcommand{\C}{\mathbb{C}}
\newcommand{\F}{\mathbb{F}}
\newcommand{\A}{\mathbb{A}}
\newcommand{\PP}{\mathbb{P}}
\newcommand{\E}{\mathbb{E}}
\newcommand{\eps}{\varepsilon}
\newcommand{\dd}{\,\mathrm{d}}
\begin{document}
\section*{UnsolvedMath ID: 30006451}
\textbf{Source problem code:} catalogue entry
\subsection{Definitions and mathematical statement}
Let \(\mathcal P=\{p\in\Z:p\ge 2\text{ and }p\text{ is prime}\}\). For every positive integer \(n\), write its unique base-ten expansion as
\[
n=\sum_{j=0}^{k-1}a_j10^j,\qquad k\ge1,\quad a_j\in\{0,1,\ldots,9\},\quad a_{k-1}\ne0.
\]
Define its digit reversal by
\[
\operatorname{rev}_{10}(n)=\sum_{j=0}^{k-1}a_j10^{k-1-j}.
\]
Leading zeros produced by reversal are discarded when interpreting the result as an integer. Call \(n\) a decimal palindrome if \(a_j=a_{k-1-j}\) for every \(0\le j<k\), equivalently if \(\operatorname{rev}_{10}(n)=n\).

The supplied record asks the following two questions, formulated as conjectures:
\begin{enumerate}
\item For every integer \(B\ge1\), there exists \(p\in\mathcal P\) such that \(p>B\) and \(\operatorname{rev}_{10}(p)\in\mathcal P\).
\item For every integer \(B\ge1\), there exists \(p\in\mathcal P\) such that \(p>B\) and \(\operatorname{rev}_{10}(p)=p\).
\end{enumerate}
The first assertion allows palindromic primes: the source imposes no requirement that the reversed prime differ from the original prime. The second assertion therefore implies the first.

The supplied record labels both questions open. Its original excerpt describes the first as seeming at least as difficult as the twin prime conjecture; this is a qualitative comparison, not a supplied mathematical reduction. The supplied dated assessment says that sieve results reach almost-prime reversals and palindromic almost-primes, while these prime assertions remain unproved. These are source qualifications, not independently verified conclusions. The report year is unresolved: the supplied citation says 2026, whereas the report designation and DOI contain 2025.
\subsection{Short English statement}
Are there infinitely many primes whose decimal digit reversals are prime, allowing the reversal to equal the original prime? Are there infinitely many decimal palindromic primes? Both are labeled open in the supplied record; this statement-only draft requires human mathematical review.
\subsection{Sources}
\begin{itemize}
\item Supplied local UnsolvedMath registry record, permanent ID \(30006451\), title \emph{Reversible Primes and Prime Palindromes}, in the Oberwolfach Reports Open Problems collection. Exact source problem code: \texttt{catalogue entry}.
\item Supplied source citation: \emph{Analytic Number Theory (2026)}, Oberwolfach Reports, DOI \texttt{10.4171/owr/2025/51}, pp.\ 2701--2764. \url{https://doi.org/10.4171/owr/2025/51}. The literature-source label is \emph{Oberwolfach Report 51/2025}; the discrepancy between 2026 and 2025 is unresolved.
\item The supplied original excerpt associates the discussion with \emph{Primes with almost prime reverse}, C\'ecile Dartyge, joint work with Jo\"el Rivat and Cathy Swaenepoel. No separate proposer is identified; the registry gives proposed year 2025.
\item Supplied literature source: \emph{Infinitude of Palindromic Almost-Prime Numbers}. \url{https://academic.oup.com/imrn/article/2024/18/12466/7729537}.
\item Literature assessment checked at \texttt{2026-08-22}: supplied dated curation, not independently verified. It labels the questions open and states that current sieve results reach almost-prime reversals and palindromic almost-primes rather than the prime assertions. The background explicitly requests independent verification of time-sensitive status and cited claims.
\item The supplied registry labels the statement \texttt{corrected\_verified}, explaining that only clear formula delimiters, notation, grammar, and typography were repaired while preserving mathematical content and scope. This is supplied editorial metadata, not independent mathematical verification.
\end{itemize}
\end{document}
